% !TEX TS-program = lualatex
\documentclass[../ctenarsky-denik.tex]{subfiles}
\begin{document}
\kniha{
  nazev=Malý princ,
  autor=Antoine de Saint-Exupéry,
  obalka=Maly-princ-1983-Frantisek-Pavlicek-Antoine-de-Saint-Exupery.jpg,
  precteno=září 2026,
  hodnoceni=5
}

\podnadpis{Děj}
Kniha začíná vzpomínkou na dětství, jak se Exupéry pokouší nakreslit hroznýše, který zrovna tráví slona. Vždy když svou kresbu ukázal někomu dospělému, tvrdil mu, že se jedná o klobouk. Dospělí mu radili aby s místo malování věnoval raději škole. Zde ukazuje, že dospělí nechápou dětský svět, protože je plný fantazie.

Příběh začíná ve chvíli, kdy pilot ztroskotá se svým letadlem uprostřed Sahary. Najednou se u něho objeví malý chlapec – malý princ. Jejich rozhovor začne tím, že malý princ chce po pilotovi, aby mu nakreslil beránka. Pilot musí beránka nakreslit několikrát znovu, protože malému princi se pořád nelíbí. Dokonce se rozčílí, když malý princ není stále spokojen, nakreslí jen bedýnku a řekne: „To je bedýnka. Beránek, kterého chceš, je uvnitř!“ Načež malý princ je konečně spokojen. Tím začne jejich přátelství.

Poté malý princ začne pilotovi vyprávět, odkud přišel, o své malé planetce B612. Ve vyprávění můžeme nepřímo nalézt různá poučení například, že Malý princ musí na planetě pravidelně plít semena baobabů. Kdyby jeden den vynechal, baobaby by se rozrostly a poté by mohly roztrhat celou jeho malou planetku.

Příběh dále pokračuje vyprávěním o květině malého prince. Na jeho planetku přiletělo neznámé semeno a vyklíčila z něj krásná rostlina. Květina se pořád pyšnila svojí krásou a malému princi poroučela, jak se o ni má starat. I když mu celou dobu poroučela, ve skutečnosti jen skrývala svoji něžnost a lásku k malému princi. On to však tehdy nechápal. Až když malý princ odjížděl, květina mu řekla, že ho má celou dobu ráda. Na své planetce měl také tři sopky, z nichž jedna byla vyhaslá. Pro jistotu ale všechny tři vymetal, aby se něco nestalo. Důležitá byla pravidelnost a zodpovědnost.

Poté co malý princ opustil svojí planetku, vydal se na cestu po ostatních planetách. Potkal zde mnoho různých lidí. Na první planetce poznal krále, který chtěl pořád někomu rozkazovat. Na celé planetě byl však sám. Z malého prince chtěl udělat svého poddaného. Král si myslel, že všechno patří jen jemu a že ho všichni poslouchají. Ve skutečnosti to bylo přesně naopak. Na druhé planetě bydlel domýšlivec. Chtěl, aby ho malý princ obdivoval. Na třetí planetě bydlel pijan. Když se ho malý princ zeptal, proč pije, odpověděl: „Abych zapomněl. – Abych zapomněl, že se stydím. – Stydím se, že piji…“ Na čtvrté planetě bydlel byznysmen. Nedělal nic jiného, než že počítal, kolik má hvězd. Myslel si, že všechno je jen jeho majetek. Nezajímalo ho nic jiného, než kolik má hvězd a co s nimi poté může udělat. Na páté planetě žil lampář se svítilnou. Pořád dokola jen rozsvěcoval a zhasínal svítilnu. Pořád dokola. Když se ho malý princ zeptal, proč si neodpočine, lampář mu odpoví, že musí vždy ráno a večer zhasínat a rozsvěcet. Je to příkaz. Jenže planeta se točí čím dál tím rychleji, a proto nemá čas na odpočinek. Šestá planeta byla obrovská. Bydlel na ní starý pán – zeměpisec. Nikdy ale nikam nejel. Malému princi řekl, že by měl jet prozkoumat Zemi.

Jeho sedmá navštívená planeta byla Země. Na poušti potkal hada. Ten mu řekl, že když někoho uštkne, může ho vrátit zpátky na zem, ze které přišel. Poté potká květinu. Ta mu řekne, kde jsou nějací lidé, malý princ se vydá tím směrem.
Šel pořád dál, až došel do zahrady plné růží. Zde byly všechny květiny jako ta, kterou měl na své planetce. Ta mu přitom tvrdila, že je jediná na světě. Byl smutný z toho, že si myslel, že má doma unikát a přitom má květinu stejnou jako všechny ostatní.
Najednou se kus od něho objevila liška. Ale nepřibližovala se k němu. Vysvětlila mu, že se k němu nemůže přiblížit, protože není ochočená. Nejdříve si musí vytvořit pouta, aby mu dovolila se k ní přiblížit. Teď pro sebe nejsou ničím důležitým. Jsou jen jedni z mnoha. Ale poté, co si vytvoří pouta, budou jen oni dva. Postupně se malý princ s liškou spřátelil. Každý den mu dovolila o trochu víc se přiblížit. Když nastal čas jeho odchodu, brečela. Ale obilí jí připomínalo malého prince, protože mělo stejnou barvu jako jeho zlaté vlasy. Při rozloučení mu řekla liška moudro: „Stáváš se navždy zodpovědným za to, cos k sobě připoutal“.

Malý princ všechny své zážitky vyprávěl pilotovi. Na poušti byli již osmý den. Došla jim skoro voda. Vydali se na cestu ke studni. Cestou si povídají. Malý princ bude mít výročí jeho příletu na Zem. Malý princ připravuje něco, o čem pilot neví. Princ ho žádá, ať jde opravovat svůj stroj. Další večer se tu sejdou. Když pilot další večer přišel, viděl, jak malý princ seděl na zdi a povídal si s hadem. Malý princ si domlouval návrat na svoji planetu. Pilot byl smutný. Poté se malý princ nechal uštknout hadem a vrátil se zpátky domů.

Příběh končí po šesti letech po této události. Pilot říká, že pokaždé, když se podívá na hvězdy, vzpomene si na malého prince.




\postavy{
\postava{Malý princ}{malý, blonďatý a kudrnatý chlapec, který se vydal cestovat do světa kvůli potížím s jeho květinou se snahou poznat smysl života. Postava je plná nevinnosti, porozumění a zvídavých otázek.}
\postava{Pilot}{rozuměl si s Malým princem, přiblížil se mu díky jeho pochopení pro dětský styl myšlení, který už ostatní dospělí dávno ztratili.}
\postava{Květina}{pyšná růže, která si myslela, že je jediná svého druhu. Nedávala znát své city.}
\postava{Liška}{moudrá postava, která se nechala ochočit Malým princem a ukázala mu veliký význam přátelství a jedinečnosti jeho květiny.}
\postava{Král}{vládne hvězdám. Každý člověk je jeho poddaný a musí ho poslouchat, protože jsou jeho rozkazy rozumné. Když se  malý princ rozhodne odejít jmenuje ho svým vyslancem.}
\postava{Domýšlivec}{chce být všemi obdivován, ale žije na své planetce sám. Slyší jen samý obdiv.}
\postava{Pijan}{pije, aby zapomněl na to, že se stydí , že pije. Je nešťastný a osamocený.}
\postava{Byznysmen}{který je neustále zaneprázdněný počítáním hvězd, o kterých si myslí, že jsou jeho. Přeje si je použít, aby mohl koupit více hvězd. Je velmi přesný a vážný.}
\postava{Lampář}{žije na planetce, kde se den a noc střídají jednou za minutu. Před nějakou dobou mu bylo nařízeno rozsvěcet lampu v noci a zhasínat ji ráno. V té době se těleso otáčelo rozumnou rychlostí a lampář měl čas na odpočinek. Poté se ale rychlost otáčení zvýšila. Lampář chtěl zůstat věrný příkazu, a tak rozsvěcí a zhasíná lampu každou minutu, bez odpočinku.}
\postava{Zeměpisec}{myslí si, že je nejdůležitější na světě. Tráví všechen svůj čas tvorbou map, ale nikdy neopouští svůj stůl k dalším objevům.}

}

\podnadpis{Ukázka}
\begin{ukazka}
  Správně vidíme jen srdcem. Oči jsou slepé.
  \zdrojukazky{kapitola XXI}
\end{ukazka}

\muj{Kniha se mi líbila, protože\ldots}
\end{document}
